% math0 manuscript. Build with: latexmk -pdf paper.tex
\documentclass[11pt]{amsart}
\usepackage[margin=1.1in]{geometry}
\usepackage{amsmath,amssymb,amsthm,mathtools}
\usepackage{url}
\usepackage{hyperref}
\usepackage{xcolor}
\usepackage{listings}
\usepackage{newunicodechar}

\lstdefinelanguage{lean}{
  keywords={theorem,lemma,def,abbrev,example,fun,by,have,show,from,let,in,if,then,else,match,with,
    namespace,end,open,import,section,variable,instance,structure,where,calc,obtain,intro,exact,
    refine,rcases,rintro,simp,omega,decide,norm_num,linarith,nlinarith,positivity,induction,cases,
    private,rw},
  sensitive=true, comment=[l]{--}, morecomment=[s]{/-}{-/}, string=[b]"}
\lstset{language=lean, basicstyle=\ttfamily\footnotesize, columns=fullflexible, keepspaces=true, breaklines=true,
  keywordstyle=\color{blue!60!black}\bfseries, commentstyle=\color{gray}, frame=single,
  extendedchars=true, inputencoding=utf8, upquote=true,
  literate={ℕ}{{$\mathbb{N}$}}1 {ℤ}{{$\mathbb{Z}$}}1 {ℝ}{{$\mathbb{R}$}}1 {∀}{{$\forall$}}1
    {∑}{{$\textstyle\sum$}}1 {∈}{{$\in$}}1 {→}{{$\to$}}1 {∧}{{$\wedge$}}1 {≤}{{$\le$}}1 {≠}{{$\ne$}}1
    {⟨}{{$\langle$}}1 {⟩}{{$\rangle$}}1}
% Unicode commonly found in Lean sources.
\newunicodechar{ℕ}{\ensuremath{\mathbb{N}}}
\newunicodechar{ℤ}{\ensuremath{\mathbb{Z}}}
\newunicodechar{ℚ}{\ensuremath{\mathbb{Q}}}
\newunicodechar{ℝ}{\ensuremath{\mathbb{R}}}
\newunicodechar{ℂ}{\ensuremath{\mathbb{C}}}
\newunicodechar{∀}{\ensuremath{\forall}}
\newunicodechar{∃}{\ensuremath{\exists}}
\newunicodechar{→}{\ensuremath{\to}}
\newunicodechar{←}{\ensuremath{\leftarrow}}
\newunicodechar{↔}{\ensuremath{\leftrightarrow}}
\newunicodechar{↦}{\ensuremath{\mapsto}}
\newunicodechar{¬}{\ensuremath{\neg}}
\newunicodechar{∧}{\ensuremath{\wedge}}
\newunicodechar{∨}{\ensuremath{\vee}}
\newunicodechar{≤}{\ensuremath{\le}}
\newunicodechar{≥}{\ensuremath{\ge}}
\newunicodechar{≠}{\ensuremath{\ne}}
\newunicodechar{∈}{\ensuremath{\in}}
\newunicodechar{∉}{\ensuremath{\notin}}
\newunicodechar{⊆}{\ensuremath{\subseteq}}
\newunicodechar{∩}{\ensuremath{\cap}}
\newunicodechar{∪}{\ensuremath{\cup}}
\newunicodechar{∑}{\ensuremath{\sum}}
\newunicodechar{∏}{\ensuremath{\prod}}
\newunicodechar{⟨}{\ensuremath{\langle}}
\newunicodechar{⟩}{\ensuremath{\rangle}}
\newunicodechar{·}{\ensuremath{\cdot}}
\newunicodechar{∘}{\ensuremath{\circ}}
\newunicodechar{⁻}{\ensuremath{^{-}}}
\newunicodechar{¹}{\ensuremath{^{1}}}
\newunicodechar{²}{\ensuremath{^{2}}}
\newunicodechar{₀}{\ensuremath{_{0}}}
\newunicodechar{₁}{\ensuremath{_{1}}}
\newunicodechar{₂}{\ensuremath{_{2}}}
\newunicodechar{α}{\ensuremath{\alpha}}
\newunicodechar{β}{\ensuremath{\beta}}
\newunicodechar{ε}{\ensuremath{\varepsilon}}
\newunicodechar{δ}{\ensuremath{\delta}}
\newunicodechar{π}{\ensuremath{\pi}}
\newunicodechar{σ}{\ensuremath{\sigma}}
\newunicodechar{φ}{\ensuremath{\varphi}}
\newunicodechar{∞}{\ensuremath{\infty}}
\newunicodechar{⌊}{\ensuremath{\lfloor}}
\newunicodechar{⌋}{\ensuremath{\rfloor}}
\newunicodechar{⌈}{\ensuremath{\lceil}}
\newunicodechar{⌉}{\ensuremath{\rceil}}
\newunicodechar{∣}{\ensuremath{\mid}}
\newunicodechar{≡}{\ensuremath{\equiv}}
\newunicodechar{×}{\ensuremath{\times}}
\newunicodechar{⧸}{\ensuremath{/}}
\newunicodechar{‖}{\ensuremath{\|}}
\newunicodechar{√}{\ensuremath{\surd}}
\newunicodechar{∅}{\ensuremath{\emptyset}}
\newunicodechar{⦃}{\ensuremath{\{\!|}}
\newunicodechar{⦄}{\ensuremath{|\!\}}}
\newunicodechar{𝓝}{\ensuremath{\mathcal{N}}}
\newunicodechar{≫}{\ensuremath{\gg}}
\newunicodechar{≪}{\ensuremath{\ll}}

\theoremstyle{plain}
\newtheorem{theorem}{Theorem}
\newtheorem{lemma}[theorem]{Lemma}
\newtheorem{proposition}[theorem]{Proposition}
\newtheorem{corollary}[theorem]{Corollary}
\theoremstyle{definition}
\newtheorem{definition}[theorem]{Definition}
\newtheorem{remark}[theorem]{Remark}

\emergencystretch=3em
% Lean identifiers in running text (breakable at dots and underscores).
\expandafter\def\expandafter\UrlBreaks\expandafter{\UrlBreaks\do\_\do\.}
\DeclareUrlCommand\decl{\urlstyle{tt}}
\newcommand{\N}{\mathbb{N}}
\newcommand{\Z}{\mathbb{Z}}
\newcommand{\R}{\mathbb{R}}
\newcommand{\up}{{\uparrow}}
\newcommand{\dn}{{\downarrow}}
\newcommand{\ev}{\operatorname{ev}}
\newcommand{\rev}{\operatorname{rev}}
\newcommand{\code}{\operatorname{code}}
\newcommand{\NC}{\mathrm{NC}}

\title[Collinear five-point subsets of $\{0,\dots,4\}^n$]{The grid $\{0,1,2,3,4\}^n$ has $(7^n-5^n)/2$ collinear five-point subsets: a formal proof of Hardin's conjecture in OEIS A081200}
\author{math0}
\address{Produced by Claude (an AI model developed by Anthropic) in the math0 project}
\date{October 24, 2026}

\begin{document}
\begin{abstract}
We prove that for every $n\ge0$ the number of $5$-element subsets of the grid $\{0,1,2,3,4\}^n\subset\R^n$ that lie on a real line is $a(n)=\sum_{k=0}^{n}\binom nk 6^{n-k}(k\bmod 2)=(7^n-5^n)/2$, where $a$ is OEIS A081200, the 6th binomial transform of $0,1,0,1,\dots$; this is the conjecture that Ron Hardin recorded in that entry in 2010.
The result is not new: these five-point sets are the winning lines of $n$-dimensional $5^n$ tic-tac-toe, and the number $((k+2)^n-k^n)/2$ of winning lines on the board of side $k$ is classical (Moser; Erd\H{o}s and Selfridge, 1973; Golomb and Hales, 2002).
Our contribution is, to our knowledge, the first formal proof, and it includes the elementary step that every real line through five grid points is such a winning line.
The theorem is verified in Lean~4 against the frozen, fidelity-reviewed statement that math0 formalized from the OEIS entry, whose definition of $a$ the kernel checks against all $23$ terms listed there.
\end{abstract}
\maketitle

\section{Introduction}

For $n\ge0$ let $G_n=\{0,1,2,3,4\}^n$, viewed as a subset of $\R^n$, and let $c(n)$ be the number of $5$-element subsets of $G_n$ that are contained in a line of $\R^n$.
The OEIS sequence A081200 \cite{A081200} (Paul Barry, 2003; offset $0$) is defined as the 6th binomial transform of $0,1,0,1,\dots$ (A000035 \cite{A000035}), that is, $a(n)=\sum_{k=0}^{n}\binom nk 6^{n-k}(k\bmod 2)$; it begins $0,1,12,109,888,6841,\dots$.
On May~24, 2010, Ron Hardin added the comment that conjecturally (``verified up to a(9)'') $a(n)$ is the number of collinear $5$-tuples of points in the $n$-dimensional $5\times\dots\times5$ grid, that is, $c(n)=a(n)$. Hardin made the analogous comments for side $4$ in A081199 \cite{A081199} and for side $6$ in A081201 \cite{A081201}. The side-$3$ analogue, that there are $(5^n-3^n)/2$ lines through $3$ points of the grid $\{0,1,2\}^n$, is recorded in A005059 \cite{A005059} by D.~W. Wilson (c.~1999), without a conjecture label. The comment in A081200 is still labelled ``Conjecture'' in the current revision (\#53, September~22, 2025). We prove it (Theorem~\ref{thm:main}) and verify the proof in Lean~4.

\subsection*{Prior work and status}
The $5$-element collinear subsets of $G_n$ are the winning lines of $n$-dimensional tic-tac-toe on the $5^n$ board. On the board of side $k$ there are $((k+2)^n-k^n)/2$ winning lines; Erd\H{o}s and Selfridge state this count in their paper on positional games \cite{ES}.
Golomb and Hales \cite[Theorem~1]{GH} give two proofs.
The first embeds the board in the board of side $k+2$: every winning line ends in exactly two cells of the border shell, and every border cell ends a line.
The second describes an oriented line by choosing, for each coordinate, ``increasing'', ``decreasing'' or one of the $k$ constants, excluding the all-constant choices; this gives $(k+2)^n-k^n$ oriented lines, and the count is halved because each line has two orientations.
Golomb and Hales state that, apart from some later sections, their work dates from before 1972.
The shell argument is credited to Leo Moser in Honsberger's \emph{Mathematical Morsels} \cite{Ho}, as reproduced in \cite{CTK} for the $8\times8\times8$ cube, which has $(10^3-8^3)/2=244$ winning lines.
The case $k=4$, $n=3$ is the game Qubic, with $(6^3-4^3)/2=76$ winning lines, which Patashnik showed by a computer-assisted search to be a first-player win \cite{Pa}; positional games on the boards $k^n$ go back to Hales and Jewett \cite{HJ}.
For $k=5$ the count is $(7^n-5^n)/2=a(n)$ (Lemma~\ref{lem:aeq}).

These sources leave one step implicit: every line of $\R^n$ containing five points of $G_n$ is a winning line in the coordinate description, that is, each coordinate of the five points, suitably ordered, is constant, runs through $0,\dots,4$, or runs through $4,\dots,0$. (Golomb and Hales define winning paths as paths ``in any straight line, including any type of diagonal''.)
This step is elementary, since each coordinate is an affine function of the position along the line (Proposition~\ref{prop:exists}).
Hardin's conjecture therefore follows from a result known since at least 1973, and we classify our contribution as the \emph{first formal proof of a known result}; the proof is the word-encoding argument, together with the step just described.
We do not know of an earlier machine-checked proof. Mathlib's formalization of the Hales--Jewett theorem \cite{Mathlib} uses \emph{combinatorial} lines, in which every coordinate is constant or increasing. For $n\ge2$ these are a proper subset of the lines counted here: for $n=2$ there are $6^2-5^2=11$ combinatorial lines, but $12$ collinear five-point sets, since the anti-diagonal is missing.
Our statement is not among those of the formal-conjectures repository \cite{FC}.

\section{Statement}\label{sec:statement}

Call a set $S\subseteq\R^n$ \emph{collinear} if there are $p,v\in\R^n$ such that every point of $S$ has the form $p+rv$ with $r\in\R$; equivalently, the linear span of $\{x-y: x,y\in S\}$ has dimension at most $1$.
Let $\mathcal C_n$ be the set of $5$-element collinear subsets of $G_n$, so that $c(n)=|\mathcal C_n|$.

\begin{theorem}\label{thm:main}
For every $n\ge0$,
\[
  |\mathcal C_n| \;=\; \sum_{k=0}^{n}\binom nk\,6^{\,n-k}\,(k\bmod 2)\;=\;\frac{7^n-5^n}{2}.
\]
\end{theorem}

The COMMENTS section of A081200 \cite{A081200} contains, verbatim:
\begin{quote}
{\small\raggedright\texttt{Conjecture (verified up to a(9)): Number of collinear 5-tuples of points in a 5 X 5 X 5 X ... n-dimensional cubic grid. - \_Ron Hardin\_, May 24 2010}\par}
\end{quote}

\subsection*{Lean formalization}
math0 formalized the statement in the file \path{lean/Math0/Statements/OEIS/A081200.lean}. It was reviewed for fidelity and frozen (SHA-256 \texttt{41ab2753\ldots 3232}) before any proof attempt. The definitions and the statement read as follows (verbatim; the import, the documentation and the data lemma \decl{a_data} are omitted):
\begin{lstlisting}
namespace Math0.Statements.A081200
def binomialTransform (m : ℕ) (b : ℕ → ℕ) (n : ℕ) : ℕ :=
  ∑ k ∈ Finset.range (n + 1), n.choose k * m ^ (n - k) * b k
def A000035 (k : ℕ) : ℕ := k % 2
def a (n : ℕ) : ℕ := binomialTransform 6 A000035 n
def toReal (n : ℕ) (x : Fin n → Fin 5) : Fin n → ℝ := fun i => ((x i : ℕ) : ℝ)
def statement : Prop :=
  ∀ n : ℕ,
    {s : Finset (Fin n → Fin 5) |
        s.card = 5 ∧ Collinear ℝ (toReal n '' (s : Set (Fin n → Fin 5)))}.ncard = a n
end Math0.Statements.A081200
\end{lstlisting}
This encodes the first equality of Theorem~\ref{thm:main}; the second equality is Lemma~\ref{lem:aeq} below, which is also proved in the Lean file (\decl{a_eq}). The encoding is as follows.
\begin{itemize}
\item \decl{a} is the binomial transform named in the OEIS entry, in the OEIS convention $(m,b)\mapsto\bigl(n\mapsto\sum_{k=0}^n\binom nk m^{n-k}b(k)\bigr)$; it is not defined by a closed form. The lemma \decl{a_data} checks all $23$ listed terms $a(0),\dots,a(22)$ (from $0$ to $1953718431395986212$) by \texttt{decide +kernel}.
\item $G_n$ is \texttt{Fin n → Fin 5}, and the injective map \decl{toReal} sends a grid point to the point of $\R^n$ (\texttt{Fin n → ℝ}) with the same coordinates. Mathlib's \texttt{Collinear ℝ S} says that the vector span of $S$ has rank at most $1$; by \decl{collinear_iff_exists_forall_eq_smul_vadd} this is our notion of collinearity.
\item A ``collinear $5$-tuple'' is read as an unordered set, a \texttt{Finset} of $5$ grid points. The data force this reading: $a(1)=1$, whereas ordered $5$-tuples would give $120$; and $a(2)=12$ counts five rows, five columns and two diagonals.
\item The count is \texttt{Set.ncard} of a set of \texttt{Finset}s of a finite type; that set is finite, so \texttt{ncard} is its true cardinality. The statement covers all $n\in\N$, matching the offset $0$; for $n=0$ both sides are $0$.
\end{itemize}
Independent brute-force counts by the math0 referees gave $c(1),\dots,c(4)=1,12,109,888=a(1),\dots,a(4)$.

\section{Proof}\label{sec:proof}

The proof follows the Lean file (Section~\ref{sec:lean}); each step names its Lean declaration. We write $[5]=\{0,1,2,3,4\}$, so that $G_n=[5]^n$.

\subsection{Words and their lines}
Let $A=\{0,1,2,3,4,\up,\dn\}$. In Lean, $A$ is \texttt{Fin 7}, with $5$ standing for $\up$ and $6$ for $\dn$.
The letters $0,\dots,4$ are \emph{constant} letters and $\up,\dn$ are \emph{arrows}.
For $c\in A$ and $t\in[5]$ define $\ev(c,t)\in[5]$ by
\[
  \ev(c,t)=c\ \ (c\in[5]),\qquad \ev(\up,t)=t,\qquad \ev(\dn,t)=4-t .
\]
Let $\rev\colon A\to A$ swap $\up$ and $\dn$ and fix the constant letters. For $u_0,u_4\in[5]$ let $\code(u_0,u_4)$ be $u_0$ if $u_0=u_4$; $\up$ if $u_0\ne u_4$ and $u_0=0$; and $\dn$ otherwise.
A \emph{word} is an element $w\in A^n$. It defines the points $p_w(t)\in G_n$ for $t\in[5]$, with coordinates $p_w(t)_i=\ev(w_i,t)$, and its \emph{line}
\[
  L(w)=\{p_w(t): t\in[5]\}\subseteq G_n .
\]
Let $\NC_n\subseteq A^n$ be the set of \emph{non-constant} words, those with at least one arrow, and write $\rev(w)$ for the word $(\rev(w_i))_i$.
In Lean (where \texttt{Fin.rev t} is $4-t$ and \texttt{rw7} is $\rev$):
\begin{lstlisting}
private def ev (c : Fin 7) (t : Fin 5) : Fin 5 :=
  if h : c.val < 5 then ⟨c.val, h⟩ else if c.val = 5 then t else Fin.rev t
private def rw7 (c : Fin 7) : Fin 7 := if c = 5 then 6 else if c = 6 then 5 else c
private def code (u0 u4 : Fin 5) : Fin 7 :=
  if u0 = u4 then ⟨u0.val, by omega⟩ else if u0.val = 0 then 5 else 6
private def pt {n : ℕ} (w : Fin n → Fin 7) (t : Fin 5) : Fin n → Fin 5 := fun i => ev (w i) t
private def line {n : ℕ} (w : Fin n → Fin 7) : Finset (Fin n → Fin 5) :=
  Finset.univ.image (pt w)
private def T : Finset (Fin 7) := Finset.univ.filter (fun c => c.val < 5)
private def NC (n : ℕ) : Finset (Fin n → Fin 7) := Finset.univ \ Fintype.piFinset (fun _ => T)
\end{lstlisting}

\subsection{Finite checks}
In Lean the next three lemmas (the first split into six declarations) are closed propositions quantified over \texttt{Fin 7} and \texttt{Fin 5}, each proved by \texttt{decide +kernel}: the kernel evaluates its \texttt{Decidable} instance by definitional unfolding to \texttt{true}. Unlike \texttt{native\_decide}, this runs no compiled code and adds no axiom. We also give proofs by hand.

\begin{lemma}[\decl{ev_inj}, \decl{ev_const}, \decl{ev_invol}, \decl{ev_rw7}, \decl{rw7_ge}, \decl{ev_affine}]\label{lem:letters}
For all $c\in A$ and $t,u\in[5]$:
\begin{enumerate}
\item if $c$ is an arrow, then $t\mapsto\ev(c,t)$ is injective, and $\ev(c,\ev(c,t))=t$;
\item if $c$ is constant, then $\ev(c,t)=\ev(c,0)$;
\item $\ev(\rev c,t)=\ev(c,4-t)$, and if $c$ is an arrow, then $\rev c$ is an arrow other than $c$;
\item $\ev(c,t)=\ev(c,0)+t\,\bigl(\ev(c,1)-\ev(c,0)\bigr)$ in $\Z$.
\end{enumerate}
\end{lemma}

\begin{proof}
The maps $t\mapsto t$ and $t\mapsto 4-t$ are involutions of $[5]$, which gives~(1); (2) and~(3) are immediate from the definitions. For~(4), the three kinds of letters give $\ev(c,\cdot)$ equal to $t\mapsto c$, $t\mapsto t$ and $t\mapsto 4-t$, all affine in $t$.
\end{proof}

\begin{lemma}[\decl{H}]\label{lem:H}
Let $u_0,u_1,u_4,t,v\in[5]$ satisfy
\[
  u_4-u_0=4(u_1-u_0)\qquad\text{and}\qquad t\,(u_4-u_0)=4(v-u_0)
\]
in $\Z$. Then $\ev(\code(u_0,u_4),t)=v$.
\end{lemma}

\begin{proof}
By the first equation $u_4-u_0$ is a multiple of $4$ in $[-4,4]$, so it is $0$, $4$ or $-4$.
If it is $0$, then $\code(u_0,u_4)=u_0$, and the second equation gives $v=u_0$.
If it is $4$, then $u_0=0$ and $u_4=4$, so $\code(u_0,u_4)=\up$, and the second equation gives $v=t$.
If it is $-4$, then $u_0=4\ne0$ and $u_4=0$, so $\code(u_0,u_4)=\dn$, and the second equation gives $v=4-t$.
In Lean the lemma is checked over all $5^5=3125$ tuples $(u_0,u_1,u_4,t,v)$.
\end{proof}

\begin{lemma}[\decl{F}]\label{lem:F}
Let $a,a'$ be arrows and $b,b'\in A$, and suppose that $\ev(b,\ev(a,\ev(a',t)))=\ev(b',t)$ for all $t\in[5]$. If $a'=a$ then $b'=b$, and if $a'\ne a$ then $b'=\rev b$.
\end{lemma}

\begin{proof}
The map $b\mapsto\ev(b,\cdot)$ from $A$ to functions $[5]\to[5]$ is injective: its values are the five distinct constant functions, the identity and $t\mapsto4-t$.
If $a'=a$, then $\ev(a,\ev(a,t))=t$ by Lemma~\ref{lem:letters}(1), so $\ev(b,\cdot)=\ev(b',\cdot)$ and $b'=b$.
If $a'\ne a$, then $\{a,a'\}=\{\up,\dn\}$ and $\ev(a,\ev(a',t))=4-t$, so $\ev(b',t)=\ev(b,4-t)=\ev(\rev b,t)$ by Lemma~\ref{lem:letters}(3), and $b'=\rev b$.
In Lean the lemma is checked over all $7^4=2401$ quadruples $(a,a',b,b')$, each with the five values of $t$.
\end{proof}

\subsection{Lines of words}

\begin{lemma}[\decl{pt_inj}, \decl{card_line}, \decl{line_collinear}]\label{lem:line}
For every word $w$, the set $L(w)$ is collinear. If $w\in\NC_n$, then $t\mapsto p_w(t)$ is injective, so $|L(w)|=5$.
\end{lemma}

\begin{proof}
Let $v\in\R^n$ have coordinates $v_i=\ev(w_i,1)-\ev(w_i,0)$. By Lemma~\ref{lem:letters}(4), $p_w(t)=p_w(0)+t\,v$ for every $t$, so $L(w)$ lies on the line through $p_w(0)$ with direction $v$.
If $w_i$ is an arrow, then $p_w(t)_i=\ev(w_i,t)$ is injective in $t$ by Lemma~\ref{lem:letters}(1).
\end{proof}

\begin{lemma}[\decl{line_rev}]\label{lem:rev}
For every word $w$, $L(\rev w)=L(w)$.
\end{lemma}

\begin{proof}
By Lemma~\ref{lem:letters}(3), $p_{\rev w}(t)=p_w(4-t)$, and $t\mapsto4-t$ is a bijection of $[5]$.
\end{proof}

\begin{lemma}[\decl{line_eq}]\label{lem:lineeq}
If $w,w'\in\NC_n$ and $L(w')=L(w)$, then $w'=w$ or $w'=\rev w$.
\end{lemma}

\begin{proof}
Choose $i_0$ such that $w_{i_0}$ is an arrow. Let $t\in[5]$. Since $p_{w'}(t)\in L(w)$, there is $u\in[5]$ with $p_w(u)=p_{w'}(t)$. Comparing coordinate $i_0$ gives $\ev(w_{i_0},u)=\ev(w'_{i_0},t)$. Applying $\ev(w_{i_0},\cdot)$, an involution by Lemma~\ref{lem:letters}(1), gives $u=\sigma(t)$, where $\sigma(t)=\ev(w_{i_0},\ev(w'_{i_0},t))$. Hence
\begin{equation}\label{eq:sigma}
  p_w(\sigma(t))=p_{w'}(t)\qquad\text{for all }t\in[5].
\end{equation}
The letter $w'_{i_0}$ is an arrow. Otherwise $\ev(w'_{i_0},\cdot)$ would be constant by Lemma~\ref{lem:letters}(2), so $\sigma(1)=\sigma(0)$, and \eqref{eq:sigma} would give $p_{w'}(1)=p_{w'}(0)$, contradicting Lemma~\ref{lem:line}.
For each coordinate $j$, \eqref{eq:sigma} says $\ev(w_j,\ev(w_{i_0},\ev(w'_{i_0},t)))=\ev(w'_j,t)$ for all $t$. Lemma~\ref{lem:F} with $(a,a',b,b')=(w_{i_0},w'_{i_0},w_j,w'_j)$ gives $w'_j=w_j$ for all $j$ if $w'_{i_0}=w_{i_0}$, and $w'_j=\rev(w_j)$ for all $j$ otherwise.
\end{proof}

\subsection{Collinear sets are lines}

\begin{lemma}[\decl{P_of_collinear}]\label{lem:P}
Let $S\subseteq G_n$ be collinear. Then for all $x,y,z\in S$ and all coordinates $i,j$,
\begin{equation}\label{eq:P}
  (x_i-z_i)(y_j-z_j)=(x_j-z_j)(y_i-z_i)\qquad\text{in }\Z .
\end{equation}
\end{lemma}

\begin{proof}
Since $z\in S$ and $S$ is collinear, there is $v\in\R^n$ with $x=z+rv$ and $y=z+r'v$ for some $r,r'\in\R$ (in Lean, \decl{collinear_iff_of_mem}). Both sides of \eqref{eq:P} equal $rr'v_iv_j$. The identity holds in $\R$ between integers, hence in $\Z$.
\end{proof}

\begin{proposition}[\decl{exists_line}]\label{prop:exists}
Every $S\in\mathcal C_n$ is of the form $L(w)$ with $w\in\NC_n$.
\end{proposition}

\begin{proof}
Let $S\subseteq G_n$ be collinear with $|S|=5$. Choose distinct $y,z\in S$ and a coordinate $i_0$ with $y_{i_0}\ne z_{i_0}$.

\emph{Step 1: the map $x\mapsto x_{i_0}$ is injective on $S$.} Suppose $x,x'\in S$, $x\ne x'$ and $x_{i_0}=x'_{i_0}$, and choose $j$ with $x_j\ne x'_j$. Apply \eqref{eq:P} to the points $(y,x',x)$ in place of $(x,y,z)$ and to the coordinates $(i_0,j)$:
\[
  (y_{i_0}-x_{i_0})(x'_j-x_j)=(y_j-x_j)(x'_{i_0}-x_{i_0})=0 .
\]
Since $x'_j\ne x_j$, we get $y_{i_0}=x_{i_0}$. The same argument with $z$ in place of $y$ gives $z_{i_0}=x_{i_0}$. Hence $y_{i_0}=z_{i_0}$, a contradiction.

\emph{Step 2: the points $q_t$.} By Step~1 the image of $S$ under $x\mapsto x_{i_0}$ has $5$ elements, so it is all of $[5]$. For each $t\in[5]$ choose $q_t\in S$ with $(q_t)_{i_0}=t$.

\emph{Step 3: the word.} Define $w\in A^n$ by $w_j=\code((q_0)_j,(q_4)_j)$. Fix $t$ and $j$, and put $u_0=(q_0)_j$, $u_1=(q_1)_j$, $u_4=(q_4)_j$ and $v=(q_t)_j$.
Applying \eqref{eq:P} to the points $(q_1,q_4,q_0)$ and the coordinates $(i_0,j)$ gives $1\cdot(u_4-u_0)=(u_1-u_0)\cdot4$. Applying it to $(q_t,q_4,q_0)$ gives $t\,(u_4-u_0)=(v-u_0)\cdot4$.
By Lemma~\ref{lem:H}, $\ev(w_j,t)=v$. Hence $p_w(t)=q_t$ for all $t$.
Moreover $w_{i_0}=\code(0,4)=\up$, so $w\in\NC_n$.

\emph{Step 4: $L(w)=S$.} We have $L(w)=\{q_t: t\in[5]\}\subseteq S$. The $q_t$ are distinct, because $(q_t)_{i_0}=t$, so $|L(w)|=5=|S|$ and $L(w)=S$.
\end{proof}

\subsection{Counting}

\begin{lemma}[\decl{card_NC}, \decl{two_mul_card}]\label{lem:count}
$|\NC_n|=7^n-5^n$, and $2\,\bigl|\{L(w): w\in\NC_n\}\bigr|=|\NC_n|$.
\end{lemma}

\begin{proof}
$\NC_n$ is the complement in $A^n$ of the $5^n$ words with only constant letters, so $|\NC_n|=7^n-5^n$.
For the second claim, partition $\NC_n$ into the fibres of $w\mapsto L(w)$. Let $w\in\NC_n$ and choose $i_0$ with $w_{i_0}$ an arrow. By Lemma~\ref{lem:letters}(3), $\rev w\in\NC_n$ and $(\rev w)_{i_0}\ne w_{i_0}$, so $\rev w\ne w$. By Lemmas~\ref{lem:rev} and~\ref{lem:lineeq}, the fibre of $L(w)$ is exactly $\{w,\rev w\}$, which has two elements. Summing over the fibres gives $|\NC_n|=2\,|\{L(w): w\in\NC_n\}|$.
\end{proof}

\begin{lemma}[\decl{a_eq}]\label{lem:aeq}
$2a(n)+5^n=7^n$ for every $n\ge0$.
\end{lemma}

\begin{proof}
By the binomial theorem, $7^n=(1+6)^n=\sum_{k=0}^n\binom nk 6^{n-k}$ and $5^n=(-1+6)^n=\sum_{k=0}^n\binom nk 6^{n-k}(-1)^k$. Since $2\,(k\bmod 2)=1-(-1)^k$, subtracting gives $7^n-5^n=\sum_{k=0}^n\binom nk6^{n-k}\cdot2\,(k\bmod2)=2a(n)$. In Lean this computation is done in $\Z$ and then cast back to $\N$.
\end{proof}

\begin{proof}[Proof of Theorem~\ref{thm:main}]
By Proposition~\ref{prop:exists} and Lemma~\ref{lem:line}, $\mathcal C_n=\{L(w): w\in\NC_n\}$ (\decl{set_eq}). By Lemma~\ref{lem:count}, $2\,|\mathcal C_n|=7^n-5^n$, and by Lemma~\ref{lem:aeq}, $2a(n)=7^n-5^n$. Hence $|\mathcal C_n|=a(n)=(7^n-5^n)/2$. In Lean the last step is done by \texttt{omega} from $2\,|\mathcal C_n|=|\NC_n|$, $|\NC_n|=7^n-5^n$ (natural subtraction) and $2a(n)+5^n=7^n$.
\end{proof}

\begin{remark}[other side lengths; not formalized]\label{rem:k}
For any side $k\ge2$ (grid $\{0,\dots,k-1\}^n$, letters $0,\dots,k-1,\up,\dn$ with $\ev(\dn,t)=k-1-t$, and $k$-element collinear subsets) the same argument gives $((k+2)^n-k^n)/2$ subsets; in Lemma~\ref{lem:H} the relation at $t=1$ forces $u_{k-1}-u_0\in\{0,\pm(k-1)\}$. As the $m$-th binomial transform of $0,1,0,1,\dots$ is $((m+1)^n-(m-1)^n)/2$, this settles, in the same reading, Hardin's comments in A081199 ($k=4$) and A081201 ($k=6$); like the case $k=5$, these follow from the classical count together with the step of Proposition~\ref{prop:exists}. Only $k=5$ is formalized here.
\end{remark}

\section{Machine verification}\label{sec:lean}

\subsection*{Files and declarations}
The proof is in the file \path{lean/Math0/Frontier/OEIS/A081200.lean} of the math0 repository, \url{https://github.com/Foundation-0/math0}. It imports the frozen statement file \path{lean/Math0/Statements/OEIS/A081200.lean} (Section~\ref{sec:statement}) and the audit commands in \path{lean/Math0/Audit.lean}.
The main declaration is the theorem \decl{Math0.Frontier.A081200.result}, whose type is \decl{statement}. It is proved by rewriting with \decl{set_eq} and \decl{Set.ncard_coe_finset} and combining \decl{two_mul_card}, \decl{card_NC} and \decl{a_eq} by \texttt{omega}; all auxiliary definitions and lemmas are \texttt{private}. The file ends with the command
\begin{lstlisting}
#math0_settles Math0.Frontier.A081200.result Math0.Statements.A081200.statement
\end{lstlisting}
The proof file contains no \texttt{sorry}, \texttt{admit}, \texttt{native\_decide}, \texttt{axiom} or \texttt{set\_option}. All finite checks (\decl{ev_inj}, \decl{ev_const}, \decl{ev_invol}, \decl{ev_rw7}, \decl{rw7_ge}, \decl{ev_affine}, \decl{H}, \decl{F}) are proved by \texttt{decide +kernel} within the default heartbeat limit.

\subsection*{Audit}
The command \texttt{\#math0\_settles}, defined in \path{lean/Math0/Audit.lean} for statements that math0 formalizes itself, fails unless (1) \decl{statement} is a \texttt{Prop}-valued definition without universe parameters; (2) \decl{statement} and \decl{result} use no axioms besides \decl{propext}, \decl{Classical.choice} and \decl{Quot.sound}, which rejects \decl{sorryAx}, \decl{Lean.ofReduceBool} (from \texttt{native\_decide}) and custom axioms; and (3) the type of \decl{result} is syntactically the statement or its negation, up to metadata, with no unfolding. Here it reports:
\begin{lstlisting}[basicstyle=\ttfamily\scriptsize]
math0 audit: 'Math0.Statements.A081200.statement' depends on axioms [Classical.choice, Quot.sound, propext]
math0 audit: 'Math0.Frontier.A081200.result' depends on axioms [Classical.choice, Quot.sound, propext]
math0 statement check: 'Math0.Frontier.A081200.result' proves 'Math0.Statements.A081200.statement'
\end{lstlisting}

\subsection*{Toolchain and reproduction}
We use Lean~4 \cite{Lean4} \texttt{v4.33.1} and Mathlib \cite{Mathlib} \texttt{v4.33.1}, both pinned in \path{lean/}. The repository also pins formal-conjectures \cite{FC} at commit \texttt{c77fdd4}, which this result does not use.
To reproduce, run \path{scripts/verify.sh} from the repository root: it fetches the Mathlib build cache, builds all math0 results with \texttt{lake build} (re-running every audit command), checks that the pinned \texttt{formal\_conjectures} and \texttt{mathlib} checkouts are unmodified, and reports the number of passed statement checks. The single file is checked by \texttt{lake env lean Math0/Frontier/OEIS/A081200.lean} in the directory \path{lean}, in under two minutes on our machine, including loading Mathlib.

{\footnotesize
\begin{thebibliography}{99}

\bibitem{CTK}
A.~Bogomolny, \emph{Tic-Tac-Toe in 3D}, Cut The Knot, \url{https://www.cut-the-knot.org/proofs/tictactoe.shtml} (excerpt from \cite{Ho}).

\bibitem{ES}
P.~Erd\H{o}s and J.~L. Selfridge, \emph{On a combinatorial game}, J. Combin. Theory Ser.~A \textbf{14} (1973), no.~3, 298--301.

\bibitem{FC}
M.~Firsching, P.~Lezeau, S.~Mercuri, M.~Z. Horv\'ath, Y.~Dillies, C.~S\"onne, E.~Wieser, F.~Zhang, T.~Hubert, B.~Ag\"uera y Arcas, and P.~Kohli, \emph{Formal Conjectures: An open and evolving benchmark for verified discovery in mathematics}, arXiv:2605.13171, 2026. Repository: \url{https://github.com/google-deepmind/formal-conjectures}.

\bibitem{GH}
S.~W. Golomb and A.~W. Hales, \emph{Hypercube tic-tac-toe}, in: R.~J. Nowakowski (ed.), \emph{More Games of No Chance}, MSRI Publications \textbf{42}, Cambridge University Press, 2002, 167--182.

\bibitem{HJ}
A.~W. Hales and R.~I. Jewett, \emph{Regularity and positional games}, Trans. Amer. Math. Soc. \textbf{106} (1963), no.~2, 222--229.

\bibitem{Ho}
R.~Honsberger, \emph{Mathematical Morsels}, Mathematical Association of America, 1978.

\bibitem{Mathlib}
The mathlib Community, \emph{The Lean mathematical library}, in: Proceedings of the 9th ACM SIGPLAN International Conference on Certified Programs and Proofs (CPP 2020), ACM, 2020, 367--381.

\bibitem{Lean4}
L.~de Moura and S.~Ullrich, \emph{The Lean 4 theorem prover and programming language}, in: Automated Deduction -- CADE 28, Lecture Notes in Comput. Sci. \textbf{12699}, Springer, 2021, 625--635.

\bibitem{Pa}
O.~Patashnik, \emph{Qubic: $4\times4\times4$ tic-tac-toe}, Math. Mag. \textbf{53} (1980), no.~4, 202--216.

\bibitem{A000035}
The OEIS Foundation Inc., \url{https://oeis.org/A000035}.

\bibitem{A005059}
The OEIS Foundation Inc., \url{https://oeis.org/A005059}; comment by D.~W. Wilson (c.~1999).

\bibitem{A081199}
The OEIS Foundation Inc., \url{https://oeis.org/A081199}; conjecture by R.~H. Hardin (May 24, 2010).

\bibitem{A081200}
The OEIS Foundation Inc., \url{https://oeis.org/A081200}, entry by P.~Barry (March 11, 2003); conjecture by R.~Hardin (May 24, 2010); revision \#53, September 22, 2025.

\bibitem{A081201}
The OEIS Foundation Inc., \url{https://oeis.org/A081201}; conjecture by R.~H. Hardin (May 23, 2010).

\end{thebibliography}
}

\end{document}
